%% init
\pdfoutput=1
\documentclass{article}

%% libraries 
\usepackage{array}
\usepackage[utf8]{inputenc}
\usepackage{listings}
\usepackage{fancyhdr}
\usepackage{lastpage}
\usepackage{float}
\usepackage{placeins}
\usepackage{url}
\usepackage{xcolor}
\usepackage{hyperref}
\usepackage{nameref}
\usepackage{enumitem}
\usepackage{adjustbox}
\usepackage{amssymb}
\usepackage{amsthm}
\usepackage{amsmath}
\usepackage{mathtools}
\usepackage{algorithm}
\usepackage{algpseudocode}
\usepackage{algcompatible}
\usepackage[numbers,sort&compress]{natbib}
\usepackage{tikz}
\usepackage{pgfplots}
\usepackage{textcomp}
\usepackage{subcaption}
\usepackage{fontawesome5}
\usepackage{tabularx}
\usepackage{multirow}
\usepackage{booktabs}
\usepackage{longtable}
\usepackage{lscape}
%% \usepackage{svg} % unused; requires shell-escape, incompatible with arxiv autotex
\usepackage[none]{hyphenat}

\graphicspath{{figures/}}
\captionsetup[figure]{font=small,labelfont=bf,labelsep=period,justification=raggedright,singlelinecheck=false}

%% doc settings
\hyphenchar\font=-1 % suppress hyphenation
\setlength\parindent{0pt} % suppress indentation
\setlength{\parskip}{0.75em} % separate paragraphs
\usepackage[margin=1.2truein]{geometry} % set page margins
\renewcommand{\tabularxcolumn}[1]{m{#1}} % left-justified
\sloppy

%% suppress section numbers
\makeatletter
\renewcommand{\@seccntformat}[1]{}
\makeatother

%% link viz
\hypersetup{
    colorlinks=true,
    linkcolor=red,
    urlcolor=red,
    citecolor=red
}

%% figs
\usetikzlibrary{
    shapes,
    arrows,
    arrows.meta,
    positioning,
    backgrounds,
    shapes.multipart,
    calc
}

\tikzset{
    block/.style={draw, fill=white, rectangle, minimum height=48pt, minimum width=48pt},
    round/.style={draw, fill=white, circle},
    sum/.style={draw, fill=white, circle, node distance=2cm},
    input/.style={coordinate},
    output/.style={coordinate},
}

\pgfplotsset{compat=newest}

%% code colors 
\definecolor{codegreen}{rgb}{0,0.6,0}
\definecolor{codegray}{rgb}{0.5,0.5,0.5}
\definecolor{codepurple}{rgb}{0.58,0,0.82}
\definecolor{backcolour}{rgb}{0.95,0.95,0.92}

\lstdefinelanguage{python}{
    % backgroundcolor=\color{backcolour},   
    % commentstyle=\color{codegreen},
    % keywordstyle=\color{blue},
    % numberstyle=\tiny\color{codegray},
    % stringstyle=\color{codepurple},
    basicstyle=\ttfamily\footnotesize,
    % breakatwhitespace=false,         
    breaklines=true,                 
    captionpos=b,                    
    % keepspaces=true,                 
    % numbers=left,                    
    % numbersep=5pt,                  
    % showspaces=false,                
    % showstringspaces=false,
    % showtabs=false,                  
    tabsize=2
}

\lstdefinelanguage{json}{
    basicstyle=\ttfamily\footnotesize,
    showstringspaces=false,
    keywordstyle=\color{blue},
    stringstyle=\color{black},
    string=[s]{"}{"}
}

\lstdefinestyle{sql}{
    language=SQL,
    basicstyle=\ttfamily\footnotesize,
    keywordstyle=\color{blue}\bfseries,
    identifierstyle=\color{black},
    stringstyle=\color{red},
    commentstyle=\color{green},
    morecomment=[l]{--},
    morecomment=[s]{/*}{*/}
}

%% math ops
\DeclareMathOperator*{\argmax}{argmax}

%% algo comments
\renewcommand{\algorithmiccomment}[1]{\texttt{//#1}}

%% page nums
\pagestyle{fancy}
\fancyhf{}
\fancyfoot[C]{Pg. \thepage \space of \pageref*{LastPage}}
\renewcommand{\headrulewidth}{0pt}

%% begin doc
\begin{document}
\title{Operating Rate Estimation in Complex Networks\\Using Capacity Frontiers}
\author{
    Nicholas Kunz$^{\textbf{*}}$ \href{mailto:nhk37@cornell.edu}{nhk37@cornell.edu}\\
    H. Oliver Gao \href{mailto:hg55@cornell.edu}{hg55@cornell.edu}\\
    {\small\textit{Corresponding Author}$^{\textbf{*}}$}\\~\\
    Systems Engineering, Cornell University\\
}
\date{\today}
\maketitle
\thispagestyle{fancy}

%% abstract
\begin{abstract}
\noindent
    Complex networks exhibit expected operating rates that span orders of magnitude across domains. A shared basis for estimating these domain-specific rates connects network capacity to typical activity. We propose the Capacity-Anchored Rate Estimator (CARE), which estimates expected operating rates using the frontier from Structural Capacity Theory (SCT), a calibration based on Extreme Value Theory (EVT), and a learned utilization function conceptually analogous to the inefficiency component in Stochastic Frontier Analysis (SFA). The frontier supplies a maximum rate estimate, while structural and dynamical features account for the utilization below it. We evaluate CARE across 25 different Network Substrate Structured Systems (N3S) in 5 broad domains. These systems couple a fixed network substrate with a dynamical process. Models trained without a given domain reliably estimate rates in the domain withheld from training. Capacity frontier anchoring reduces log-rate Mean Absolute Error (MAE) by nearly 80\% when compared to the models without it. Adjusting for network capacity contracts the regression target prior to training and reduces the variability that the utilization model needs to explain. Component ablations assess the positive contributions of this decomposition. These findings support the practical utility of capacity frontiers when estimating expected operating rates across otherwise unrelated systems.
\end{abstract}

%% classifications and keywords
\begin{quotation}
\noindent\small
\textbf{Keywords:} complex networks $\mid$ capacity frontier $\mid$ dynamical processes $\mid$ graph invariants $\mid$ domain transfer
\end{quotation}

%% untitled intro
\vspace{1em}
\phantomsection
\label{sec:intro}

Expected operating rates describe how much activity complex networks sustain over an observation period. Communication throughput depends on network structure and routing, while human activity can alternate between bursts and inactivity \citep{ref35,ref113}. A system's empirical maximum therefore does not determine the rate at which it typically operates. Estimating typical activity requires accounting for both its overall scale and its separation from observed extremes. Whether an estimated capacity frontier can provide a useful reference for this task across otherwise unrelated domains is the question considered here.

Previous work provides several ways to relate system organization to aggregate activity. Resource-distribution networks offer a basis for biological scaling, while urban scaling relates infrastructure and socioeconomic activity to population size \citep{ref17,ref18}. Deviations from urban scaling separate shared size-dependent variation from characteristics specific to individual cities \citep{ref121}. Related work on network dynamics separates the contributions of topology and dynamical interactions to system response and resilience \citep{ref15,ref16}. These studies motivate representations that distinguish common constraints from system-specific behavior. They do not, by themselves, establish whether a capacity frontier learned across domains supplies a useful scale for estimating expected operating rates in a domain withheld from training.

We propose the Capacity-Anchored Rate Estimator (CARE) to estimate expected operating rates without specifying or simulating the governing dynamics of each system. CARE uses the empirical frontier of Structural Capacity Theory (SCT), which represents network structure through graph invariants and dynamical processes through process signatures \citep{ref122}. It therefore uses observed information about the dynamics without requiring an explicit mechanistic model of their evolution. Its scope is Network Substrate Structured Systems (N3S): systems with an observable process attributed to a fixed network substrate and a measurable maximum rate. We extend the estimation task from that maximum to the expected rate $\bar{r} = M/L$, where $M$ is the total event count across $L$ observed intervals. The frontier supplies a system-specific reference from which a downstream model estimates the remaining departure.

We combine the estimated SCT frontier with an extreme-value adjustment and a learned utilization correction. Extreme Value Theory (EVT) relates maxima to the distribution and number of underlying observations \citep{ref123}; here, an EVT-motivated calibration uses process variability and observation length to adjust the frontier toward typical activity. Structural and dynamical descriptors then inform the utilization correction. Separating an attainable reference from observed performance has precedents in Stochastic Frontier Analysis (SFA) and Data Envelopment Analysis (DEA) \citep{ref19,ref21}. CARE uses this distinction for rate estimation, without identifying utilization slack as economic inefficiency. The estimated frontier enters as a fixed offset on the log-rate scale, shifting the downstream task from estimating the full rate scale to estimating a correction relative to capacity. This formulation tests whether capacity and utilization provide a transferable basis for estimating activity across systems with different underlying mechanisms.

We evaluate CARE across 25 systems spanning 5 broad domains: Earth and physical sciences, life sciences and medicine, technology and information, trade and institutions, and transportation and infrastructure. These systems represent disparate phenomena whose expected operating rates span approximately 6 orders of magnitude. By withholding entire domains from training, we test whether capacity anchoring supports estimation across this heterogeneity without specifying each system's governing dynamics. We evaluate Poisson, Poisson--Tweedie, and Negative Binomial models, with discipline and random holdouts providing complementary checks. Frontier anchoring reduced held-out-domain log-rate MAE by 75--78\%, and the 2 overdispersed families each placed all 25 systems within a factor of 3 of observed rates. Component ablations, constant-ceiling and remapping controls, and a feature-matched direct-rate alternative assessed the contributions of calibration, utilization, and system-specific frontier information. Separating time-series inputs and targets provided an additional robustness check. Together, these evaluations tested whether empirical capacity estimates have practical value for estimating activity below the frontier.

%% results
\section{Results}

\subsection{Collapses Prediction Error}

Prediction error was substantially reduced while estimation-observation consensus improved. SCT-FAR was compared with unanchored rate estimation across 3 Generalized Linear Model (GLM) families and 4 resampling regimes (Fig.~2a--c; Table~2). The corpus contained 25 systems across 5 broad domains. Prediction error was measured using Mean Absolute Error (MAE) on log-transformed operating rates. Estimation-observation consensus was measured using the Consensus Index (CI), as defined in the companion study. Root Mean Squared Error (RMSE) and Geometric Mean Multiplicative Error (GMME) characterized larger errors and multiplicative disagreement, respectively.

Under domain-level Leave-One-Group-Out Cross-Validation (LOGO-CV), MAE decreased from 3.255 to 0.804 for Poisson, from 1.461 to 0.316 for Poisson--Tweedie, and from 1.450 to 0.316 for Negative Binomial estimation. Estimation-observation consensus increased from 0.602 to 0.974, from 0.823 to 0.984, and from 0.847 to 0.983, respectively. For Negative Binomial estimation, RMSE decreased from 1.954 to 0.382 and GMME from 4.715 to 1.468. The proportion of estimates within a factor of 2 of the observed rate increased from 44\% to 76\%, while the proportion within a factor of 3 increased from 56\% to 100\% (Fig.~4). Poisson--Tweedie estimation achieved the same proportions using SCT-FAR.

\subsection{Generalizes Across Domains}

Domain generalization was evident despite wide differences in subject and scale. Under domain-level LOGO-CV, the estimator was trained on 4 of the 5 broad domains and evaluated on the 5 individual disciplines in the withheld domain. Under discipline-level LOGO-CV, it was trained on 24 of the 25 disciplines and evaluated on the one that was withheld. These evaluations were compared with 5-Fold-CV and 10-Fold-CV to determine whether generalization depended on the resampling regime (Table~3).

Across the 4 regimes, Negative Binomial MAE ranged from 0.263 to 0.347 using SCT-FAR, compared with 1.444 to 1.558 under unanchored estimation. Estimation-observation consensus ranged from 0.980 to 0.983, compared with 0.803 to 0.872. RMSE ranged from 0.330 to 0.444, and GMME from 1.390 to 1.509. Poisson--Tweedie estimation exhibited comparable consistency, with MAE of 0.273--0.338 and estimation-observation consensus of 0.984--0.985.

Prediction quality was also compared across resampling regimes using random $k$-Fold-CV repeated 30 times (Fig.~2d; Supplementary Table~1). Paired Wilcoxon Two One-Sided Tests (TOST) across 15 estimator-by-domain comparisons supported equivalence within the empirically derived margins for MAE, CI, RMSE, and GMME. All Holm-adjusted probabilities were below 0.012. Thus, withholding broad domains preserved comparable aggregate performance within the specified margins.

\subsection{Relieves Scale Burden}

The target's statistical variability substantially contracted. The original log-rate target was compared with the utilization slack target obtained through frontier anchoring and EVT adjustment (Fig.~3a). Frontier values were estimated for held-out systems, and target dispersion was evaluated before fitting the operating-rate estimators. Under domain-level LOGO-CV, the Interquartile Range (IQR) decreased from 4.46 to 0.92, a contraction of 79.4\%. Variance decreased from 11.11 to 1.67, a contraction of 85.0\%.

Target contraction persisted across all 4 resampling regimes (Supplementary Table~5). IQR contraction ranged from 75.3\% to 91.2\%, while variance contraction ranged from 81.5\% to 90.4\%. Smaller fitted coefficient magnitudes accompanied these changes (Fig.~3b; Supplementary Table~6). Under Negative Binomial domain-level LOGO-CV, mean absolute coefficients across folds decreased by 74.0--94.1\% for the 3 structural features and 58.3--80.7\% for the 3 process features. These reductions accompanied the improvements in MAE, CI, RMSE, and GMME. Frontier anchoring therefore reduced the variation that the operating-rate estimators needed to explain.

\subsection{Refines Capacity-to-Rate Conversion}

Prediction quality was jointly improved by EVT calibration and GLM utilization. Their contributions were evaluated by removing the EVT adjustment, replacing utilization features with a fitted intercept, and removing both components (Fig.~5a; Table~5). Under Negative Binomial domain-level LOGO-CV, MAE increased from 0.316 to 0.393, 0.599, and 1.074, respectively. Estimation-observation consensus decreased from 0.983 to 0.982, 0.974, and 0.941. The corresponding RMSE values increased from 0.382 to 0.457, 0.727, and 1.185, while GMME increased from 1.468 to 1.602, 1.988, and 3.275.

Paired bootstrap 95\% intervals for all 3 MAE increases excluded zero, and paired Wilcoxon tests supported inferior performance after Holm adjustment (all adjusted $p < 0.001$). Retaining both components produced the lowest MAE across all 3 GLM families and 4 resampling regimes (Table~6). Under Negative Binomial 5-Fold-CV, however, removing EVT improved RMSE from 0.444 to 0.406 and estimation-observation consensus from 0.980 to 0.985, while MAE increased from 0.347 to 0.350. The combined formulation consistently improved MAE, although the ordering was not uniform across every metric.

\subsection{Substantiates Decomposition}

All tested component-wise permutations, information-matched models, arbitrary ceilings, and shuffled frontiers were inferior. Inferiority was evaluated through paired MAE comparisons with SCT-FAR under Negative Binomial domain-level LOGO-CV, with the remaining metrics and resampling regimes characterizing its extent (Fig.~5a; Tables~5--6). Beyond component ablation, alternative formulations tested whether the predictive advantage depended on the decomposition and the correspondence between systems and their frontier estimates.

An information-matched direct rate estimate increased MAE from 0.316 to 1.415 and reduced estimation-observation consensus from 0.983 to 0.827. Replacing system-specific frontiers with a common arbitrary ceiling produced MAE of 1.377 and estimation-observation consensus of 0.880. Shuffling frontiers across systems produced the largest deterioration, with corresponding values of 2.955 and 0.607. Across these controls, RMSE increased from 0.382 to 1.867--3.829 and GMME from 1.468 to 4.382--26.918.

All 7 primary MAE comparisons, including component ablations, supported inferiority after Holm adjustment (all adjusted $p < 0.001$). The information-matched, arbitrary-ceiling, and shuffled-frontier controls performed worse than SCT-FAR on all 4 metrics across all 3 families and 4 resampling regimes. These results supported the predictive contribution of the capacity-to-rate decomposition and its system-specific frontier estimates.

\subsection{Preserves Performance Temporally}

Prediction quality was generally preserved when operating rate estimation relied on samples from prior time horizons. Each series was divided chronologically at 3 proportions. The first 25\%, 50\%, or 75\% supplied the process signatures, EVT statistics, and frontier training maxima, while the remaining samples supplied the target operating rate (Fig.~5b; Supplementary Tables~2--4). Graph invariants spanned both windows.

Under the 50\% split and domain-level LOGO-CV, Negative Binomial MAE decreased from 1.515 under unanchored estimation to 0.581 using SCT-FAR. Estimation-observation consensus increased from 0.809 to 0.954. RMSE decreased from 1.976 to 0.812, and GMME from 5.005 to 1.892. Paired bootstrap intervals favored SCT-FAR on all 4 metrics across all 3 families, and paired MAE comparisons remained significant after Holm adjustment. Lower MAE, RMSE, and GMME and higher estimation-observation consensus persisted across all 36 combinations of window proportion, GLM family, and resampling regime.

Carrying the earlier-window mean forward yielded MAE of 0.474, estimation-observation consensus of 0.964, RMSE of 0.776, and GMME of 1.692 at the 50\% split. This reference achieved better aggregate scores than SCT-FAR. Nevertheless, separating the process samples from the target preserved SCT-FAR's predictive advantage over the tested unanchored GLMs across every evaluated configuration.

%% discussion
\section*{Discussion}
\label{sec:discussion}

SCT-FAR estimates expected operating rates across disciplinary boundaries. Withholding entire domains preserved prediction quality across systems with substantial differences in subject, size, scale, units, and dynamics. Information from one set of systems was therefore useful for estimating activity in another without fitting a separate model of each system's dynamical process. These findings support an approximately transferable relationship between empirical capacity and typical operation across the tested domains. Rate estimation can therefore draw on evidence beyond a system's own substantive boundary.

An empirical upper-envelope can provide an informative reference for central activity. The capacity frontier was constructed to characterize maxima, yet anchoring to it substantially improved estimates of expected operating rates. Target contraction and smaller fitted coefficient magnitudes suggest why this transfer was effective. The frontier accounted for much of the variation across systems, reducing the variation that the operating-rate estimators needed to explain. Capacity estimation thus supplied a shared reference scale for systems whose observed activity differed by several orders of magnitude. This extends the computational role of Structural Capacity Theory from describing empirical limits to estimating realized operation.

The organization of information contributed to prediction quality. Information-matched direct rate estimation performed worse than SCT-FAR, while arbitrary ceilings and shuffled frontiers failed to reproduce its performance. These comparisons supported the contribution of the capacity-to-rate decomposition and the correspondence between each system and its frontier estimate. EVT calibration and GLM utilization further improved the conversion to typical activity. Together, the results suggest that a capacity reference can make available structural and process information more effective for estimation. The frontier's predictive contribution therefore extends beyond its numerical role as an upper bound.

The separation between the frontier and observed activity can be accommodated without resolving its full physical meaning. In the companion study, substantial slack left the interpretation of the frontier's absolute level open. SCT-FAR demonstrates that such a frontier can support accurate estimation after adjustment for process variation and learned utilization. The correction may absorb measurement differences, observation-window effects, and frontier estimation error alongside differences in realized operation. Its predictive success supports the decomposition, while identifying these contributions separately remains an open question.

All findings were conditioned on the 25 tested systems, their shared representation, the selected frontier estimator, and the 3 GLM families. The advantage over unanchored estimation persisted when earlier process samples were used to estimate later operating rates, although graph invariants spanned both windows and the carry-forward mean achieved better aggregate performance at the primary split. Independent corpora, alternative representations, and different temporal scales provide direct tests of generality. Within the evaluated scope, the principal conclusion is that empirical capacity provides a transferable basis for estimating typical activity across otherwise unrelated systems.

%% methods
\section*{Methods}
\label{sec:methods}

%% implementation
\subsection*{Implementation and Reproducibility}
\label{sub-sec:imp-rep}

All analyses were implemented in \texttt{Python 3.11.9} \citep{ref86,ref102}. All dependency versions were pinned in the repository's \texttt{requirements.txt}. Graphs, processes, and encodings used \texttt{python-igraph 1.0.0} \citep{ref48,ref93}, \texttt{NumPy 2.3.5} \citep{ref81,ref96}, \texttt{SciPy 1.16.3} \citep{ref82,ref97}, and \texttt{pandas 3.0.0} \citep{ref83,ref84}. Urban street network construction used \texttt{networkx 3.5} \citep{ref91,ref98} and \texttt{OSMnx 2.0.6} \citep{ref118,ref119}.

Public dataset retrieval was conducted using \texttt{requests 2.32.5} \citep{ref120}. Temporal graph datasets were loaded using \texttt{torch-geometric 2.7.0} \citep{ref89,ref99} and \texttt{torch\_geometric\_temporal 0.56.2} \citep{ref90,ref100}. Observational records for gravitational wave and meteorological systems were retrieved using \texttt{gwosc 0.8.1} and \texttt{meteostat 1.7.6}, respectively \citep{ref116,ref117}.

Estimators used \texttt{scikit-learn 1.8.0} \citep{ref72,ref94}, \texttt{sklearn-quantile 0.1.1} \citep{ref87}, \texttt{XGBoost 3.1.1} \citep{ref39,ref95}, and \texttt{PyTorch 2.9.0} \citep{ref85,ref115}. Distance correlation was computed using \texttt{dcor 0.7} \citep{ref88}, and scientific figures were generated with \texttt{matplotlib 3.9.2} \citep{ref92,ref101}. Unspecified settings followed library defaults.

All stochastic procedures started from random seed 42 with repeat indices applied as fixed numerical offsets. Random permutations, cross-validation folds, and perturbation realizations originated from seeded generators. 

\texttt{PyTorch 2.9.0} deterministic algorithms were executed in warn-only mode and execution defaulted to CUDA when GPU resources were available and CPU otherwise. Results reproduced with the code repository, corpus, and environment specifications are subject to floating-point imprecision.
\clearpage

\section*{Data, Materials, and Software Availability}
\label{sec:data}
\label{sec:code}
The corpus was collated from publicly available data sources (\textit{SI Appendix}, Table~\ref{tab:si-corpus}). The processed records, scripts for data processing, frontier estimation, evaluation, perturbation, falsification, and ablation, along with the notebooks that reproduce every figure and table, are distributed with the source code repository at \url{https://github.com/nickkunz/network-capacity}.

\section*{Acknowledgments}
\label{sec:acknowledge}
We would like to thank F. Parise, W. Xu, N. de Santi, and S. Visvesvaran for their early feedback, discussions, and encouragement.

\section*{Author Contributions}
\label{sec:authors}
N. Kunz initiated and developed the concept, methods, implementation, testing, documentation, and manuscript. H. O. Gao initiated and supervised the development of the investigation and manuscript. The authors have agreed to the submitted version of this manuscript.

\section*{Funding}
\label{sec:funding}
This work was supported by Cornell University and Gao Labs.

\section*{Competing Interests}
\label{sec:competing}
The authors declare no competing interests.

\section*{Additional Information}
\label{sec:additional}
\textbf{Supporting information} is in \textit{SI Appendix}, Tables~\ref{tab:si-corpus}--\ref{tab:si-protocol} and \hyperref[sec:si-note-bipartite]{Notes~S1}--\hyperref[sec:si-note-processing]{S2}.

\textbf{Correspondence and requests for materials} should be addressed to N. Kunz.

%% references
\newpage
\bibliographystyle{unsrtnat}
\bibliography{ref}

\clearpage
\begingroup
\renewcommand{\tablename}{Table}
\renewcommand{\thetable}{S\arabic{table}}
\renewcommand{\theHtable}{supp.\arabic{table}}
\setcounter{table}{0}
\captionsetup[table]{font=footnotesize,labelfont=bf,labelsep=period,justification=raggedright,singlelinecheck=false}
\setlength{\tabcolsep}{3pt}
\renewcommand{\arraystretch}{1.15}
\footnotesize

\section*{SI Appendix}
\begingroup
\captionsetup{hypcap=false}
\captionof{table}{Corpus composition, definitions, sources, and date provenance.}
\label{tab:si-corpus}
\endgroup

\begin{center}
\rotatebox{90}{%
\begin{minipage}{.92\textheight}
\scriptsize
\setlength{\tabcolsep}{2pt}
\begin{tabular}{@{}
    >{\raggedright\arraybackslash}p{\dimexpr.08\linewidth-2\tabcolsep\relax}
    >{\raggedright\arraybackslash}p{\dimexpr.10\linewidth-2\tabcolsep\relax}
    >{\raggedright\arraybackslash}p{\dimexpr.10\linewidth-2\tabcolsep\relax}
    >{\raggedright\arraybackslash}p{\dimexpr.10\linewidth-2\tabcolsep\relax}
    >{\raggedright\arraybackslash}p{\dimexpr.10\linewidth-2\tabcolsep\relax}
    >{\raggedright\arraybackslash}p{\dimexpr.10\linewidth-2\tabcolsep\relax}
    >{\raggedright\arraybackslash}p{\dimexpr.10\linewidth-2\tabcolsep\relax}
    >{\raggedright\arraybackslash}p{\dimexpr.08\linewidth-2\tabcolsep\relax}
    >{\raggedright\arraybackslash}p{\dimexpr.105\linewidth-2\tabcolsep\relax}
    >{\raggedright\arraybackslash}p{\dimexpr.135\linewidth-2\tabcolsep\relax}@{}}
\toprule
\textbf{System} & \textbf{Domain} & \textbf{Discipline} & \textbf{Nodes} & \textbf{Edges} & \textbf{Event} & \textbf{Target} & \textbf{Time (Days)} & \textbf{Source} & \textbf{Notes} \\
\midrule
\textbf{\textit{gwosc}} & Earth \& Physical Sciences & Cosmology &
Gravitational-\allowbreak wave detectors (H1, L1, V1, K1, G1). & Complete graph among detectors. &
GWTC event records from any GWTC catalog attributed to at least one corpus detector, including the marginal and auxiliary catalogs. & Maximum daily records $r_{\max} = 3$, $y^* = 1.386$. & 2015-09-14 to 2024-01-09 ($n = 3{,}040$). &
GWOSC event API (\ref{src:S1}, \ref{src:S24}, \ref{src:S25}, \ref{src:S26}, \ref{src:S27}, \ref{src:S30}). & Catalog articles published in 2019, 2021, 2023, and 2024. GWTC-4.0 first released 26 August 2025 and its results article was published in 2026. Fixed event cutoff 2024-01-09. \\
\addlinespace[3.5pt]
\textbf{\textit{river}} & Earth \& Physical Sciences & Hydrology &
USGS stream gauges on the Colorado River (HUC 15010001, 15010002, 15010005). & Non-crossing path among gauges. &
Peaks of discharge (parameter 00060) above the site-wise 99th percentile. & Maximum daily discharge-peak exceedances $r_{\max} = 9$, $y^* = 2.303$. & 2014-11-12 to 2023-08-01 ($n = 3{,}185$). &
USGS NWIS site and instantaneous-\allowbreak values services (\ref{src:S2}). & NWIS dataset record published 1994. API query cutoff 2024-12-31. \\
\addlinespace[3.5pt]
\textbf{\textit{auger}} & Earth \& Physical Sciences & Physics &
Pierre Auger Observatory surface-\allowbreak detector stations. & Reciprocal 6-nearest neighbours within $1.1\times$ array spacing (1,500 m and 750 m arrays). &
Surface-\allowbreak detector air-shower events. & Maximum daily air-shower events $r_{\max} = 66$, $y^* = 4.205$. & 2004-01-03 to 2018-09-01 ($n = 5{,}356$). &
Auger Open Data station map and Zenodo record 10488964 (\ref{src:S3}). & Open Data version 3 released 2024. Observations span 2004--2018. \\
\addlinespace[3.5pt]
\textbf{\textit{seismic}} & Earth \& Physical Sciences & Geology &
IRIS IU network stations. & Complete graph among stations. &
Reviewed USGS events, 2023. & Maximum daily reviewed events $r_{\max} = 1{,}129$, $y^* = 7.030$. & 2023-01-01 to 2023-12-30 ($n = 364$). &
IRIS FDSN station service (\ref{src:S4}) and USGS FDSN event service (\ref{src:S5}). & GSN dataset record published 1988. ComCat record published 2017. Event-query cutoff 2023-12-31. \\
\addlinespace[3.5pt]
\textbf{\textit{rain}} & Earth \& Physical Sciences & Meteorology &
Meteostat hourly stations in Laos. & Delaunay triangulation of station coordinates. &
Station-hours with precipitation $>1$ mm. & Maximum daily wet station-hours $r_{\max} = 46$, $y^* = 3.850$. & 2024-01-12 to 2024-03-30 ($n = 79$). &
Meteostat (\ref{src:S6}). & Mutable service accessed 1 September 2026. Query cutoff 2024-03-31. \\
\addlinespace[3.5pt]
\textbf{\textit{celegans}} & Life Sciences \& Medicine & Biology &
\textit{C. elegans} developmental stages. & Known stage-to-stage transitions (normal, dauer, and recovery pathways), entered undirected. &
Developmental transitions at published mean times. & Maximum daily developmental transitions $r_{\max} = 8$, $y^* = 2.197$. & Days 0--33 ($n = 34$). &
Cassada \& Russell, \textit{Dev. Biol.} \textbf{46}, 326--342 (1975) (\ref{src:S7}), and Golden \& Riddle, \textit{Dev. Biol.} \textbf{102}, 368--378 (1984) (\ref{src:S8}). & Timing studies published 1975 and 1984. Days are developmental-age indices, not calendar dates. \\
\midrule
\multicolumn{10}{r@{}}{\textit{Continued on next page}}\\
\end{tabular}
\end{minipage}}
\end{center}

\begin{landscape}
\begingroup
\scriptsize
\setlength{\tabcolsep}{2pt}
\addtocounter{table}{-1}
\renewcommand{\theHtable}{supp.1.cont}
\begin{longtable}{@{}
    >{\raggedright\arraybackslash}p{\dimexpr.08\linewidth-2\tabcolsep\relax}
    >{\raggedright\arraybackslash}p{\dimexpr.10\linewidth-2\tabcolsep\relax}
    >{\raggedright\arraybackslash}p{\dimexpr.10\linewidth-2\tabcolsep\relax}
    >{\raggedright\arraybackslash}p{\dimexpr.10\linewidth-2\tabcolsep\relax}
    >{\raggedright\arraybackslash}p{\dimexpr.10\linewidth-2\tabcolsep\relax}
    >{\raggedright\arraybackslash}p{\dimexpr.10\linewidth-2\tabcolsep\relax}
    >{\raggedright\arraybackslash}p{\dimexpr.10\linewidth-2\tabcolsep\relax}
    >{\raggedright\arraybackslash}p{\dimexpr.08\linewidth-2\tabcolsep\relax}
    >{\raggedright\arraybackslash}p{\dimexpr.105\linewidth-2\tabcolsep\relax}
    >{\raggedright\arraybackslash}p{\dimexpr.135\linewidth-2\tabcolsep\relax}@{}}
\multicolumn{10}{@{}l@{}}{\textbf{Table~\ref{tab:si-corpus}.} Corpus composition and provenance \textit{(continued).}}\\[4pt]
\toprule
\textbf{System} & \textbf{Domain} & \textbf{Discipline} & \textbf{Nodes} & \textbf{Edges} & \textbf{Event} & \textbf{Target} & \textbf{Time (Days)} & \textbf{Source} & \textbf{Notes} \\
\midrule
\endfirsthead
\multicolumn{10}{@{}l@{}}{\textbf{Table~\ref{tab:si-corpus}.} Corpus composition and provenance \textit{(continued).}}\\[4pt]
\toprule
\textbf{System} & \textbf{Domain} & \textbf{Discipline} & \textbf{Nodes} & \textbf{Edges} & \textbf{Event} & \textbf{Target} & \textbf{Time (Days)} & \textbf{Source} & \textbf{Notes} \\
\midrule
\endhead
\midrule
\multicolumn{10}{r@{}}{\textit{Continued on next page}}\\
\endfoot
\bottomrule
\endlastfoot
\textbf{\textit{crop}} & Life Sciences \& Medicine & Ecology &
Crops and pollinator taxa. & Observed crop--pollinator interactions. &
Daily flower-visitor census visits from primary study records. & Maximum daily census visits $r_{\max} = 1{,}266$, $y^* = 7.144$. & 2004-05-07 to 2012-05-29 ($n = 2{,}945$). &
CropPol repository (Rader et al. primary census records) (\ref{src:S9}, \ref{src:S31}). & Database article published 2022. Observations span 2004--2012. \\
\addlinespace[5pt]
\textbf{\textit{chickenpox}} & Life Sciences \& Medicine & Epidemiology &
Hungarian counties. & County adjacency from the Chickenpox dataset. &
Weekly county case counts summed across counties, divided by 7 to a daily average. & Maximum daily-average cases $r_{\max} = 378$, $y^* = 5.938$. & 2005-01-03 to 2014-12-29 ($n = 3{,}648$). &
Rozemberczki et al. Chickenpox dataset and UCI Hungarian chickenpox cases (\ref{src:S10}). & UCI dataset released 2021. Observations span 2005--2014. Source reports weekly totals, converted to daily averages. \\
\addlinespace[5pt]
\textbf{\textit{faers}} & Life Sciences \& Medicine & Pharmacology &
Medicinal products and MedDRA reaction terms co-reported with imatinib. & Drug--reaction co-occurrence. &
Unique FAERS case receipts. & Maximum daily case receipts $r_{\max} = 45$, $y^* = 3.829$. & 2004-01-23 to 2025-12-31 ($n = 8{,}014$). &
openFDA drug-event API, query IMATINIB (\ref{src:S11}). & Mutable API retrieved 13 March 2026. Case-receipt cutoff 2025-12-31. \\
\addlinespace[5pt]
\textbf{\textit{epilepsy}} & Life Sciences \& Medicine & Neurology &
19 scalp electrodes (10--20 system). & Anatomical neighbours on the 10--20 montage. &
Annotated seizures in CHB-MIT (chb01--chb24). & Maximum daily annotated seizures $r_{\max} = 27$, $y^* = 3.332$. & Relative recording days 0--76 ($n = 77$). &
PhysioNet CHB-MIT (\ref{src:S12}). & PhysioNet version 1.0.0 released 2010. Calendar dates are unavailable, so recording order is expressed as relative days. \\
\addlinespace[5pt]
\textbf{\textit{federal}} & Trade \& Institutions & Government &
Awarding agencies and recipients. & Agency--\allowbreak recipient contract pairs. &
Waterfowl-\allowbreak related contract start dates. & Maximum daily contract starts $r_{\max} = 17$, $y^* = 2.890$. & 2011-01-19 to 2024-12-30 ($n = 5{,}095$). &
USAspending API, keyword waterfowl (\ref{src:S13}). & Mutable API retrieved 25 August 2026. Award-start cutoff 2024-12-31. \\
\addlinespace[5pt]
\textbf{\textit{mooc}} & Trade \& Institutions & Education &
Students and course actions. & Complete bipartite graph between students and actions. &
Student--action timestamps. & Maximum daily student actions $r_{\max} = 18{,}213$, $y^* = 9.810$. & Days 0--29 ($n = 30$). &
SNAP act-mooc (\ref{src:S14}). & Dataset article published 2019. Timestamps reindexed to relative days. \\
\addlinespace[5pt]
\textbf{\textit{amazon}} & Trade \& Institutions & Commerce &
Reviewers and book products. & Reviewer--\allowbreak product review pairs. &
Review timestamps. & Maximum daily reviews $r_{\max} = 37{,}293$, $y^* = 10.527$. & 1996-05-20 to 2014-07-23 ($n = 6{,}639$). &
SNAP Amazon product graph, Books ratings (\ref{src:S15}). & Dataset article published 2015. Reviews span 1996--2014. \\
\addlinespace[5pt]
\textbf{\textit{bitcoin}} & Trade \& Institutions & Cryptocurrency &
OTC users and integer rating values. & User--\allowbreak rating--\allowbreak user incidences over the complete Bitcoin OTC history. &
Rating timestamps. & Maximum daily ratings $r_{\max} = 319$, $y^* = 5.768$. & 2010-11-08 to 2016-01-25 ($n = 1{,}905$). &
PyTorch Geometric BitcoinOTC (\ref{src:S16}). & Dataset article published 2016. Ratings span 2010--2016. \\
\addlinespace[5pt]
\textbf{\textit{world}} & Trade \& Institutions & International Development &
World Bank lending instruments and countries. & Instrument--\allowbreak country eligibility. &
Project board-approval dates. & Maximum daily project approvals $r_{\max} = 12$, $y^* = 2.565$. & 2014-01-14 to 2024-12-20 ($n = 3{,}994$). &
World Bank projects and country APIs (\ref{src:S17}). & Mutable APIs retrieved 29 September 2026. Approval-year cutoff 2024. \\
\addlinespace[5pt]
\textbf{\textit{wiki}} & Technology \& Information & Knowledge Graphs &
Wikipedia mathematics pages. & Page--page relations from WikiMaths snapshots. &
Daily summed page activity. & Maximum daily summed page activity $r_{\max} = 1{,}323{,}743$, $y^* = 14.096$. & 2019-03-16 to 2021-03-15 ($n = 731$). &
UCI Wikipedia Math Essentials and PyTorch Geometric Temporal WikiMaths (\ref{src:S18}). & UCI dataset released 2021. Activity spans 2019--2021. \\
\addlinespace[5pt]
\textbf{\textit{jodie}} & Technology \& Information & Interaction Networks &
Wikipedia users and pages. & User--page interactions. &
Interaction timestamps. & Maximum daily interactions $r_{\max} = 6{,}087$, $y^* = 8.714$. & Days 0--30 ($n = 31$). &
PyTorch Geometric JODIE, Wikipedia (\ref{src:S14}). & Dataset article published 2019. Timestamps reindexed to relative days. \\
\addlinespace[5pt]
\textbf{\textit{overflow}} & Technology \& Information & Online Forums &
Stack Overflow users and calendar days. & Complete bipartite graph between users and days. &
Comment, answer, and other timed interactions. & Maximum daily timed interactions $r_{\max} = 52{,}446$, $y^* = 10.868$. & 2008-08-01 to 2016-03-06 ($n = 2{,}775$). &
SNAP sx-stackoverflow (\ref{src:S19}). & Dataset article published 2017. Interactions span 2008--2016. \\
\addlinespace[5pt]
\textbf{\textit{email}} & Technology \& Information & Email Exchanges &
European research-\allowbreak institution email accounts and days. & Complete bipartite graph between accounts and days. &
Email timestamps. & Maximum daily emails $r_{\max} = 2{,}270$, $y^* = 7.728$. & Days 0--803 ($n = 804$). &
SNAP email-\allowbreak Eu-\allowbreak core-\allowbreak temporal (\ref{src:S19}). & Dataset article published 2017. Timestamps reindexed to relative days. \\
\addlinespace[5pt]
\textbf{\textit{college}} & Technology \& Information & Social Networks &
UC Irvine college-message users and days. & Complete bipartite graph between users and days. &
Private-message timestamps. & Maximum daily private messages $r_{\max} = 2{,}678$, $y^* = 7.893$. & 2004-04-15 to 2004-10-26 ($n = 195$). &
SNAP CollegeMsg (\ref{src:S20}). & Dataset article published 2009. Messages were recorded in 2004. \\
\addlinespace[5pt]
\textbf{\textit{idling}} & Transportation \& Infrastructure & Mobility &
Halifax drive-network nodes from the unsimplified OSM graph. & Street segments (OSM drive graph, undirected). &
Vehicle-idling records. & Maximum daily idling records $r_{\max} = 732$, $y^* = 6.597$. & 2025-06-02 to 2025-09-02 ($n = 93$). &
OpenStreetMap via OSMnx, query Halifax, Canada (\ref{src:S21}, \ref{src:S22}), and local idling CSVs (\ref{src:S29}). & Data preprint first posted 2024 and revised 2025. Observations are from 2025. OSMnx article published 2025. \\
\addlinespace[5pt]
\textbf{\textit{windmill}} & Transportation \& Infrastructure & Energy &
Wind turbines. & Complete graph among turbines. &
Raw-production on/off transitions (threshold $10^{-6}$). & Maximum daily turbine transitions $r_{\max} = 7$, $y^* = 2.079$. & Days 0--727 ($n = 728$). &
PyTorch Geometric Temporal Windmill\allowbreak Output\allowbreak Large \citep{ref90}. & Loader article published 2021. Source supplies snapshot order rather than calendar dates. \\
\addlinespace[5pt]
\textbf{\textit{metrla}} & Transportation \& Infrastructure & Traffic &
METR-LA loop detectors. & Detector adjacency from METR-LA snapshots. &
Onset of congestion (raw speed at or below the detector-wise 1st percentile). & Maximum daily congestion onsets $r_{\max} = 675$, $y^* = 6.516$. & 2012-03-01 to 2012-06-27 ($n = 119$). &
PyTorch Geometric Temporal METRLA \citep{ref90} (\ref{src:S23}). & Benchmark article published 2018 and loader article 2021. Sensor observations are from 2012. \\
\addlinespace[5pt]
\textbf{\textit{pemsbay}} & Transportation \& Infrastructure & Traffic II &
PEMS-BAY loop detectors. & Detector adjacency from PEMS-BAY snapshots. &
Onset of congestion (raw speed at or below the detector-wise 1st percentile). & Maximum daily congestion onsets $r_{\max} = 634$, $y^* = 6.454$. & 2017-01-01 to 2017-06-29 ($n = 180$). &
PyTorch Geometric Temporal PEMS-BAY \citep{ref90} (\ref{src:S23}). & Benchmark article published 2018 and loader article 2021. Sensor observations are from 2017. \\
\addlinespace[5pt]
\textbf{\textit{montevideo}} & Transportation \& Infrastructure & Ridership &
Montevideo bus stops. & Stop adjacency from MontevideoBus snapshots. &
Stop-hours with standardized inflow above the global 99th percentile. & Maximum daily exceedance stop-hours $r_{\max} = 256$, $y^* = 5.549$. & Days 0--29 ($n = 30$). &
PyTorch Geometric Temporal MontevideoBus \citep{ref90} (\ref{src:S28}). & Data recorded in October 2020. Loader article published 2021. Dates represented as relative days. \\
\end{longtable}
\par\smallskip
\begin{minipage}{\linewidth}
\footnotesize
Each row represents one N3S, with observations aggregated to $r_t$ (Eq.~\eqref{eq:series}) and targets $r_{\max}$ and $y^* = \ln(1 + r_{\max})$ (Eq.~\eqref{eq:maximum}). Metrla, pemsbay, montevideo, windmill, and wiki provide complete daily series, including zero-event days. Chickenpox provides weekly county totals, summed and divided by 7, with one daily-average value per report date; the other 24 systems use daily event counts. The remaining 19 series include only days with recorded events. Windows follow native recording periods, expressed as dates or relative day indices; $n$ is the inclusive span in days, not necessarily the series length $L$. The 7 API systems (gwosc, river, seismic, rain, faers, federal, and world) also use fixed script query dates, with end cutoffs in the source notes. Durations were not standardized, and maxima remain conditional on their windows. Observation dates differ from source release or publication dates; source notes specify versions and relative-date conventions. Additional processing conventions are provided in the supplement (\hyperref[sec:si-note-processing]{Note~S2}). Citations S1--S31 refer to the SI sources (\nameref{sec:si-sources}).
\end{minipage}
\endgroup
\end{landscape}

\clearpage
\begin{table}[H]
\centering
\footnotesize
\captionsetup{position=top,skip=12pt}
\caption{Domain transfer and baseline frontier performance.}
\label{tab:si-transfer}
\begin{tabularx}{\linewidth}{@{}lXrrrr@{}}
\toprule
Evaluation & Method & Median EI [Q1--Q3] & VR & MV & MS \\
\midrule
Transfer & Domain (LOGO-CV) & 0.75 [0.74--0.77] & 0.04 & 0.08 & 6.86 \\
 & Discipline (LOGO-CV) & 0.73 [0.72--0.76] & 0.00 & 0.00 & 6.12 \\
\addlinespace
Baseline & Random ($5$-Fold-CV) & 0.77 [0.76--0.78] & 0.02 & 0.06 & 6.86 \\
 & Random ($10$-Fold-CV) & 0.76 [0.75--0.78] & 0.01 & 0.02 & 6.87 \\
\bottomrule
\end{tabularx}
\par\vspace{12pt}
\begin{minipage}{\linewidth}
\footnotesize
Entries are medians of estimator-level means over folds or groups and 30 repeats (seeds 42--71), with quartiles across 9 estimators in brackets. Groups, folds, repeats, and estimators are equally weighted. LOGO-CV folds are fixed across repeats; random folds are reshuffled. EI denotes the Efficiency Index (\nameref{sec:performance-metrics},~Eq.~\eqref{eq:ei}), VR the violation rate, MV the mean exceedance among within-fold violators (0 if none), and MS the mean non-negative slack across evaluated systems. MV and MS use natural-log target units. Descriptive summaries only.
\end{minipage}
\end{table}

\begin{table}[H]
\centering
\footnotesize
\captionsetup{position=top,skip=12pt}
\caption{Equivalence of resampling regimes to domain-level LOGO-CV.}
\label{tab:si-transfer-equivalence}
\begin{tabularx}{\linewidth}{@{}Xrrrrrr@{}}
\toprule
\multicolumn{7}{c}{\textbf{Domain Transfer Equivalence}} \\
\cmidrule(lr){1-7}
Evaluation & \shortstack{Median\\$\Delta\mathrm{EI}$} & \shortstack{Rank-biserial\\$r$} & \shortstack{TOST\\$p$} & \shortstack{Holm-adj.\\$p$} & Sig. & Eq. \\
\midrule
Discipline (LOGO-CV) & 0.00 & -0.08 & $<0.001$ & $<0.001$ & *** & Yes \\
Random ($5$-Fold-CV) & 0.00 & 0.40 & $<0.001$ & $<0.001$ & *** & Yes \\
Random ($10$-Fold-CV) & 0.00 & 0.24 & $<0.001$ & $<0.001$ & *** & Yes \\
\midrule
\multicolumn{7}{c}{\textbf{Estimation-Observation Consensus Equivalence}} \\
\cmidrule(lr){1-7}
Evaluation & \shortstack{Median\\$\Delta\mathrm{CI}$} & \shortstack{Rank-biserial\\$r$} & \shortstack{TOST\\$p$} & \shortstack{Holm-adj.\\$p$} & Sig. & Eq. \\
\midrule
Discipline (LOGO-CV) & 0.00 & -0.14 & $<0.001$ & $<0.001$ & *** & Yes \\
Random ($5$-Fold-CV) & -0.01 & -0.32 & $<0.001$ & $<0.001$ & *** & Yes \\
Random ($10$-Fold-CV) & 0.00 & -0.34 & $<0.001$ & $<0.001$ & *** & Yes \\
\bottomrule
\end{tabularx}
\par\vspace{12pt}
\begin{minipage}{\linewidth}
\footnotesize
Each contrast uses 45 paired estimator-by-domain cells (9 estimators $\times$ 5 domains). Out-of-sample predictions were averaged over 30 seeds (42--71) before computing EI and CI within the same 5 broad domains for every regime. EI and CI are defined in the methodology (\nameref{sec:performance-metrics},~Eqs.~\eqref{eq:ei} and~\eqref{eq:ci}). Shifts are medians of paired differences (test comparison $-$ domain-level LOGO-CV), not differences of marginal medians. Rank-biserial $r$ uses nonzero, unshifted differences with average ranks for ties; positive values favor the comparison regime. TOST uses paired one-sided Wilcoxon signed-rank tests on $\Delta+\delta$ and $\Delta-\delta$, assuming a symmetric paired-difference distribution with location $\theta$. The nulls are $\theta\leq-\delta$ and $\theta\geq\delta$; unadjusted $p$ is the larger of the two one-sided values. Margins $\delta_{\mathrm{EI}} = 0.02$ and $\delta_{\mathrm{CI}} = 0.11$ are IQRs across the 9 reference estimator-level means after equal averaging across domains, floored to 2 decimals. Raw differences with magnitude $<10^{-10}$ were zeroed before shifting. Tests used \texttt{SciPy 1.16.3}, \texttt{zero\_method='wilcox'}, \texttt{method='auto'}, and \texttt{correction=False}, discarding zeros after shifting. $p_{\mathrm H}$ is Holm-adjusted across the 3 regime comparisons separately for EI and CI. Eq. requires rejection of both nulls at $p_{\mathrm H}<0.05$. Shared estimators, training systems, and domains make inference conditional on this dependent panel (\nameref{sec:si-stat-conventions}).
\end{minipage}
\end{table}

\clearpage
\begin{table}[H]
\centering
\footnotesize
\captionsetup{position=top,skip=12pt}
\caption{Perturbation effects under the highest test intensity.}
\label{tab:si-perturb}
\setlength{\tabcolsep}{2pt}
\begin{tabularx}{\linewidth}{@{}l>{\raggedright\arraybackslash}Xrrrrrr@{}}
\toprule
\multicolumn{8}{c}{\textbf{Domain Transfer}} \\
\cmidrule(lr){1-8}
Perturbation & Method & \shortstack{Median\\$\Delta\mathrm{EI}$} & \shortstack{Rank-biserial\\$r$} & \shortstack{TOST\\$p$} & \shortstack{Holm-adj.\\$p$} & Sig. & Eq. \\
\midrule
Network & Sample & 0.00 & -0.35 & $<0.001$ & $<0.001$ & *** & Yes \\
 & Rewire & 0.00 & -0.23 & $<0.001$ & $<0.001$ & *** & Yes \\
 & Densify & 0.00 & -0.40 & $<0.001$ & $<0.001$ & *** & Yes \\
\addlinespace
Invariants & Jitter & 0.00 & -0.45 & $<0.001$ & $<0.001$ & *** & Yes \\
 & Noise & 0.00 & -0.60 & $<0.001$ & $<0.001$ & *** & Yes \\
 & Subset & 0.00 & 0.05 & $<0.001$ & $<0.001$ & *** & Yes \\
\addlinespace
Process & Bootstrapping & 0.01 & 0.33 & $<0.001$ & $<0.001$ & *** & Yes \\
 & Scaling & -0.01 & -0.45 & $<0.001$ & $<0.001$ & *** & Yes \\
 & Smoothing & -0.02 & -0.75 & $<0.001$ & $<0.001$ & *** & Yes \\
\addlinespace
Signatures & Jitter & 0.00 & -0.31 & $<0.001$ & $<0.001$ & *** & Yes \\
 & Noise & 0.00 & -0.44 & $<0.001$ & $<0.001$ & *** & Yes \\
 & Subset & -0.01 & -0.61 & $<0.001$ & $<0.001$ & *** & Yes \\
\midrule
\multicolumn{8}{c}{\textbf{Estimation-Observation Consensus}} \\
\cmidrule(lr){1-8}
Perturbation & Method & \shortstack{Median\\$\Delta\mathrm{CI}$} & \shortstack{Rank-biserial\\$r$} & \shortstack{TOST\\$p$} & \shortstack{Holm-adj.\\$p$} & Sig. & Eq. \\
\midrule
Network & Sample & 0.00 & -0.58 & $<0.001$ & $<0.001$ & *** & Yes \\
 & Rewire & 0.00 & -0.43 & $<0.001$ & $<0.001$ & *** & Yes \\
 & Densify & 0.00 & -0.47 & $<0.001$ & $<0.001$ & *** & Yes \\
\addlinespace
Invariants & Jitter & 0.00 & -0.44 & $<0.001$ & $<0.001$ & *** & Yes \\
 & Noise & 0.00 & -0.42 & $<0.001$ & $<0.001$ & *** & Yes \\
 & Subset & 0.00 & -0.41 & $<0.001$ & $<0.001$ & *** & Yes \\
\addlinespace
Process & Bootstrapping & -0.11 & -0.82 & 0.028 & 0.028 & * & Yes \\
 & Scaling & -0.02 & -0.23 & $<0.001$ & $<0.001$ & *** & Yes \\
 & Smoothing & -0.06 & -0.73 & 0.002 & 0.003 & ** & Yes \\
\addlinespace
Signatures & Jitter & -0.01 & -0.29 & $<0.001$ & $<0.001$ & *** & Yes \\
 & Noise & -0.01 & -0.68 & $<0.001$ & $<0.001$ & *** & Yes \\
 & Subset & -0.04 & -0.82 & $<0.001$ & $<0.001$ & *** & Yes \\
\bottomrule
\end{tabularx}
\par\vspace{12pt}
\begin{minipage}{\linewidth}
\footnotesize
Frozen domain-level LOGO-CV retains original-data training fits and uses 45 paired estimator-by-domain cells (9 estimators $\times$ 5 domains), summarized over 30 seeds (42--71). EI is averaged over repetitions in both conditions. CI is computed from seed-averaged predictions for the original and deterministic conditions, but averaged over paired seed--realization metrics for stochastic perturbations. EI and CI are defined in the methodology (\nameref{sec:performance-metrics},~Eqs.~\eqref{eq:ei} and~\eqref{eq:ci}). Thus aggregation order can contribute to stochastic CI shifts. Shifts are medians of paired differences (perturbed $-$ original). Rank-biserial $r$ uses nonzero, unshifted differences with average ranks for ties; positive values favor perturbation. Only each method's most severe tested setting is included (Table~\ref{tab:si-protocol}); Sample denotes node removal and coordinate subsetting. Original maximum targets are retained. TOST uses paired one-sided Wilcoxon signed-rank tests on $\Delta+\delta$ and $\Delta-\delta$, assuming a symmetric paired-difference distribution with location $\theta$. The nulls are $\theta\leq-\delta$ and $\theta\geq\delta$; unadjusted $p$ is the larger of the two one-sided values. Margins $\delta_{\mathrm{EI}} = 0.08$ and $\delta_{\mathrm{CI}} = 0.14$ are IQRs across the 45 original-condition cells, floored to 2 decimals. Raw differences with magnitude $<10^{-10}$ were zeroed before shifting. Tests used \texttt{SciPy 1.16.3}, \texttt{zero\_method='wilcox'}, \texttt{method='auto'}, and \texttt{correction=False}, discarding zeros after shifting. $p_{\mathrm H}$ is Holm-adjusted across the 12 methods separately for EI and CI. Eq. requires rejection of both nulls at $p_{\mathrm H}<0.05$. Shared estimators, training systems, and domains make inference conditional on this dependent panel (\nameref{sec:si-stat-conventions}).
\end{minipage}
\end{table}

\clearpage
\begin{table}[H]
\centering
\footnotesize
\captionsetup{position=top,skip=12pt}
\caption{Falsification effects under frozen and retrained frontiers.}
\label{tab:si-falsify}
\setlength{\tabcolsep}{2pt}
\begin{tabularx}{\linewidth}{@{}l>{\raggedright\arraybackslash}Xrrrrrr@{}}
\toprule
\multicolumn{8}{c}{\textbf{Domain Transfer}} \\
\cmidrule(lr){1-8}
Falsification & Method & \shortstack{Median\\$\Delta\mathrm{EI}$} & \shortstack{Rank-biserial\\$r$} & \shortstack{One-sided\\$p$} & \shortstack{Holm-adj.\\$p$} & Sig. & Diff. \\
\midrule
Frozen & Target Permutation & -0.02 & -0.38 & 0.013 & 0.013 & * & Yes \\
 & Random Generation & -0.04 & -0.91 & $<0.001$ & $<0.001$ & *** & Yes \\
 & Vector Generation & -0.03 & -0.74 & $<0.001$ & $<0.001$ & *** & Yes \\
\addlinespace
Retrain & Target Permutation & -0.09 & -0.53 & $<0.001$ & 0.001 & ** & Yes \\
 & Random Generation & -0.02 & -0.59 & $<0.001$ & $<0.001$ & *** & Yes \\
 & Vector Generation & -0.03 & -0.57 & $<0.001$ & $<0.001$ & *** & Yes \\
\midrule
\multicolumn{8}{c}{\textbf{Estimation-Observation Consensus}} \\
\cmidrule(lr){1-8}
Falsification & Method & \shortstack{Median\\$\Delta\mathrm{CI}$} & \shortstack{Rank-biserial\\$r$} & \shortstack{One-sided\\$p$} & \shortstack{Holm-adj.\\$p$} & Sig. & Diff. \\
\midrule
Frozen & Target Permutation & -0.29 & -0.97 & $<0.001$ & $<0.001$ & *** & Yes \\
 & Random Generation & -0.28 & -0.91 & $<0.001$ & $<0.001$ & *** & Yes \\
 & Vector Generation & -0.21 & -0.88 & $<0.001$ & $<0.001$ & *** & Yes \\
\addlinespace
Retrain & Target Permutation & -0.24 & -0.99 & $<0.001$ & $<0.001$ & *** & Yes \\
 & Random Generation & -0.18 & -0.94 & $<0.001$ & $<0.001$ & *** & Yes \\
 & Vector Generation & -0.31 & -0.97 & $<0.001$ & $<0.001$ & *** & Yes \\
\bottomrule
\end{tabularx}
\par\vspace{12pt}
\begin{minipage}{\linewidth}
\footnotesize
Each domain-level LOGO-CV contrast uses 45 paired estimator-by-domain cells (9 estimators $\times$ 5 domains). Predictions are averaged over 30 model-fitting seeds (42--71) before computing EI and CI within each domain. EI and CI are defined in the methodology (\nameref{sec:performance-metrics},~Eqs.~\eqref{eq:ei} and~\eqref{eq:ci}). Each falsified corpus was generated once with seed 42; model-fitting repeats are not independent negative-control corpora. Frozen frontiers retain original-data training fits, while retrained frontiers refit standardizers and estimators on the falsified training folds. Shifts are medians of paired differences (falsified $-$ original), with negative values favoring the original condition. Rank-biserial $r$ uses nonzero, unshifted differences with average ranks for ties. One-sided Wilcoxon signed-rank tests assess $H_0:\theta\geq0$ versus $H_1:\theta<0$, assuming a symmetric paired-difference distribution with location $\theta$. Raw differences with magnitude $<10^{-10}$ were zeroed. Tests used \texttt{SciPy 1.16.3}, \texttt{zero\_method='wilcox'}, \texttt{method='auto'}, and \texttt{correction=False}, discarding zero differences. $p$ is unadjusted; $p_{\mathrm H}$ is Holm-adjusted across the 6 method-by-frontier comparisons separately for EI and CI. Diff. requires $p_{\mathrm H}<0.05$ and a negative median shift, supporting inferiority to the original condition. Shared estimators, training systems, and domains make inference conditional on this dependent panel (\nameref{sec:si-stat-conventions}).
\end{minipage}
\end{table}

\clearpage
\begin{table}[H]
\centering
\footnotesize
\captionsetup{position=top,skip=12pt}
\caption{Non-inferiority of the log-additive decomposition to alternative formulations.}
\label{tab:si-ablation}
\setlength{\tabcolsep}{2pt}
\begin{tabularx}{\linewidth}{@{}l>{\raggedright\arraybackslash}Xrrrrrr@{}}
\toprule
\multicolumn{8}{c}{\textbf{Domain Transfer}} \\
\cmidrule(lr){1-8}
Ablation & Method & \shortstack{Median\\$\Delta\mathrm{EI}$} & \shortstack{Rank-biserial\\$r$} & \shortstack{One-sided\\$p$} & \shortstack{Holm-adj.\\$p$} & Sig. & NI. \\
\midrule
Simple & Invariants Only & -0.01 & -0.77 & $<0.001$ & $<0.001$ & *** & Yes \\
 & Signatures Only & -0.01 & -0.21 & $<0.001$ & $<0.001$ & *** & Yes \\
\addlinespace
Complex & Joint Features & -0.02 & -0.55 & $<0.001$ & $<0.001$ & *** & Yes \\
 & Interaction Terms & -0.01 & -0.46 & $<0.001$ & $<0.001$ & *** & Yes \\
\midrule
\multicolumn{8}{c}{\textbf{Estimation-Observation Consensus}} \\
\cmidrule(lr){1-8}
Ablation & Method & \shortstack{Median\\$\Delta\mathrm{CI}$} & \shortstack{Rank-biserial\\$r$} & \shortstack{One-sided\\$p$} & \shortstack{Holm-adj.\\$p$} & Sig. & NI. \\
\midrule
Simple & Invariants Only & -0.20 & -0.86 & $<0.001$ & $<0.001$ & *** & Yes \\
 & Signatures Only & -0.08 & -0.55 & $<0.001$ & $<0.001$ & *** & Yes \\
\addlinespace
Complex & Joint Features & -0.27 & -0.67 & $<0.001$ & $<0.001$ & *** & Yes \\
 & Interaction Terms & -0.10 & -0.65 & $<0.001$ & $<0.001$ & *** & Yes \\
\bottomrule
\end{tabularx}
\par\vspace{12pt}
\begin{minipage}{\linewidth}
\footnotesize
Each retrained domain-level LOGO-CV contrast uses 45 paired estimator-by-domain cells (9 estimators $\times$ 5 domains) and 30 model-fitting seeds (42--71), with all pipelines refit within their training folds. EI is averaged over repetitions; CI is computed from seed-averaged out-of-sample predictions within each domain. EI and CI are defined in the methodology (\nameref{sec:performance-metrics},~Eqs.~\eqref{eq:ei} and~\eqref{eq:ci}). Shifts are medians of paired differences (alternative formulation $-$ log-additive decomposition); negative values favor the decomposition. Rank-biserial $r$ uses nonzero, unshifted differences with average ranks for ties, independently of the tested margin. One-sided Wilcoxon signed-rank tests assess $H_0:\theta\geq\delta$ versus $H_1:\theta<\delta$, assuming a symmetric paired-difference distribution with location $\theta$. Margins $\delta_{\mathrm{EI}} = 0.08$ and $\delta_{\mathrm{CI}} = 0.14$ are IQRs across the 45 log-additive reference cells, floored to 2 decimals. Raw differences with magnitude $<10^{-10}$ were zeroed before subtracting $\delta$. Tests used \texttt{SciPy 1.16.3}, \texttt{zero\_method='wilcox'}, \texttt{method='auto'}, and \texttt{correction=False}, discarding zeros after shifting. $p$ is unadjusted; $p_{\mathrm H}$ is Holm-adjusted separately within the 2 simple and 2 complex alternatives for EI and CI. NI. requires $p_{\mathrm H}<0.05$ and a median shift below the margin. This supports non-inferiority within the stated margin, not superiority over the alternatives. Shared estimators, training systems, and domains make inference conditional on this dependent panel (\nameref{sec:si-stat-conventions}).
\end{minipage}
\end{table}

\clearpage
\begin{table}[H]
\centering
\footnotesize
\captionsetup{position=top,skip=12pt}
\caption{Graph invariant definitions, expressions, and asymptotic time complexity.}
\label{tab:si-features}
\begin{tabularx}{\linewidth}{@{}ll>{\raggedright\arraybackslash}Xll@{}}
\toprule
\textbf{Category} & \textbf{Invariant} & \textbf{Definition} & \textbf{Expression} & \textbf{Complexity} \\
\midrule
Simple & 1. Number of Nodes & Number of nodes\newline & $n = |V|$ & $O(1)$ \\
 & 2. Number of Edges & Number of edges\newline & $m = |E|$ & $O(1)$ \\
 & 3. Articulation Points & Vertices whose removal splits $G$ & $\lvert \{ v : \omega(G - v) > \omega(G) \} \rvert$ & $O(n+m)$ \\
 & 4. Bridges & Edges whose removal splits $G$ & $\lvert \{ e : \omega(G - e) > \omega(G) \} \rvert$ & $O(n+m)$ \\
\addlinespace
Cohesion & 5. Diameter & Longest shortest-path length \dag & $\max_v \epsilon(v)$ & $O(nm)$ \\
 & 6. Radius & Smallest vertex eccentricity \dag & $\min_v \epsilon(v)$ & $O(nm)$ \\
 & 7. Degeneracy & Largest $k$ with\newline a $k$-core & $k^{*} = \max_v c(v)$ & $O(n+m)$ \\
 & 8. $k$-Core Size & Vertices in the\newline core at $k^{*}$ & $\lvert \{ v : c(v) = k^{*} \} \rvert$ & $O(n+m)$ \\
\addlinespace
Extremal & 9. Maximum Degree & Highest vertex degree & $\max_v d_v$ & $O(n)$ \\
\addlinespace
Statistical & 10. Degree Variance & Spread of vertex\newline degrees & $n^{-1} \sum_v (d_v - \bar{d})^2$ & $O(n)$ \\
 & 11. Degree Entropy & Entropy of degree distribution & $-\sum_k p_k \ln p_k$ & $O(n)$ \\
 & 12. Joint Degree Entropy & Entropy of edge degree pairs & $-\sum_{k \le l} q_{kl} \ln q_{kl}$ & $O(m)$ \\
 & 13. Degree Skewness & Asymmetry of degree distribution & $\frac{\sqrt{n(n-1)}}{n-2}\, g_1$ & $O(n)$ \\
 & 14. Degree Kurtosis & Tail weight of degree distribution & $\frac{(n-1)[(n+1) g_2 + 6]}{(n-2)(n-3)}$ & $O(n)$ \\
 & 15. Global Clustering Coefficient & Fraction of\newline closed triples & $3 N_{\triangle} / N_{\wedge}$ & $O(n d_{\max}^2)$ \\
 & 16. Assortativity & Correlation of endpoint degrees & $\mathrm{corr}_{\{u,v\} \in E}(d_u, d_v)$ & $O(m)$ \\
\addlinespace
Spectral & 17. Normalized Laplacian 2nd Moment & Diffusion intensity & $\mathrm{tr}(L^2)/n$ & $O(m)$ \\
 & 18. Normalized Laplacian 3rd Moment & Diffusion asymmetry & $\mathrm{tr}(L^3)/n$ & $O(m^2)$ \\
 & 19. Random Walk 3rd Moment & Three-step recurrence & $\mathrm{tr}(P^3)/n$ & $O(m^2)$ \\
 & 20. Random Walk 4th Moment & Four-step recurrence & $\mathrm{tr}(P^4)/n$ & $O(m^2)$ \\
 & 21. Adjacency 4th Moment & Four-length walk density & $\mathrm{tr}(A^4)/n$ & $O(m^2)$ \\
\bottomrule
\end{tabularx}
\par\vspace{12pt}
\begin{minipage}{\linewidth}
\footnotesize
Admission required global graph descriptors defined solely by topology, independent of embeddings, weights, external parameters, and arbitrary thresholds, exactly and deterministically computable at corpus scale, and non-trivially variable across systems. All returned values were finite. \dag~Diameter and radius use the largest connected component; degeneracy and $k$-core size use the full graph. Spectral moments use exact trace identities without eigendecomposition; the 4 complete bipartite substrates use closed-form invariants. Here $n = |V|$, $m = |E|$, $d_v$ is vertex degree, $\bar{d}$ its mean, $\mu_j = n^{-1} \sum_v (d_v - \bar{d})^j$, $g_1 = \mu_3 / \mu_2^{3/2}$, and $g_2 = \mu_4 / \mu_2^{2} - 3$. $\omega$ counts connected components, $\epsilon(v)$ is eccentricity in hops, $c(v)$ is core number, $p_k$ is the fraction of nodes of degree $k$, $q_{kl}$ is the fraction of edges with unordered endpoint degrees $\{k, l\}$, and $N_{\triangle}$ and $N_{\wedge}$ count triangles and connected triples. $A$ is the adjacency matrix, $\hat{A} = D^{-1/2} A D^{-1/2}$ over non-isolated nodes, $L = I - \hat{A}$, and $P = D^{-1} A$. Isolated nodes count toward $n$ but contribute 0 to $\mathrm{tr}(L^k)$. Logarithms are natural; assortativity counts both edge orientations. Undefined assortativity and clustering are 0. Degree skewness and kurtosis are 0 for equal degrees or $n < 3$ and $n < 4$, respectively. References: \citep{ref49,ref50,ref51,ref52,ref61,ref62,ref63,ref64,ref79,ref112,ref114}.
\end{minipage}
\end{table}

\clearpage
\begin{table}[H]
\centering
\footnotesize
\captionsetup{position=top,skip=12pt}
\caption{Process signature definitions, expressions, and asymptotic time complexity.}
\label{tab:si-signatures}
\begin{tabularx}{\linewidth}{@{}ll>{\raggedright\arraybackslash}Xll@{}}
\toprule
\textbf{Property} & \textbf{Signature} & \textbf{Definition} & \textbf{Expression} & \textbf{Complexity} \\
\midrule
Dependence & 1. Lag-1 Autocorrelation & Short-term sequential dependence & $\frac{\sum_{i=1}^{L-1} \tilde{r}_i \tilde{r}_{i+1}}{\sum_{i=1}^{L} \tilde{r}_i^{2}}$ & $O(L)$ \\
 & 2. Hurst Exponent & Long-range\newline dependence & $\ln(R/s) / \ln L$ & $O(L)$ \\
\addlinespace
Dispersion & 3. Coefficient of Variation & Scale-free\newline dispersion & $s / \bar{r}$ & $O(L)$ \\
 & 4. Fano Factor & Deviation from Poisson dispersion & $s^{2} / \bar{r}$ & $O(L)$ \\
\addlinespace
Volatility & 5. Normalized Successive Difference & Local temporal volatility & $\frac{1}{(L-1)\,\bar{r}} \sum_{i=1}^{L-1} \lvert r_{i+1} - r_i \rvert$ & $O(L)$ \\
\addlinespace
Trend & 6. Trend Coefficient & Systematic non-stationarity & $\mathrm{corr}(r_i, i)$ & $O(L)$ \\
\addlinespace
Asymmetry & 7. Count Skewness & Distributional asymmetry & $\frac{1}{L} \sum_{i=1}^{L} (\tilde{r}_i / s)^{3}$ & $O(L)$ \\
\bottomrule
\end{tabularx}
\par\vspace{12pt}
\begin{minipage}{\linewidth}
\footnotesize
Admission required global series descriptors defined solely by the rate series, without external parameters or arbitrary thresholds, exactly and deterministically computable, and non-trivially variable across systems. All returned values were finite. For the chronologically ordered rate series $r_1, \ldots, r_L$ (Eq.~\eqref{eq:series}), $L = \lvert T \rvert$, $\bar{r}$ is the mean, $\tilde{r}_i = r_i - \bar{r}$, and $s$ is the sample standard deviation with denominator $L - 1$. Interval inclusion follows corpus definitions (Table~\ref{tab:si-corpus}); lag-1 dependence and trend index consecutive included observations, not calendar gaps. The Hurst exponent is a full-series, single-window rescaled-range estimate clipped to $[0, 1]$, with $R = \max_i Y_i - \min_i Y_i$ and cumulative sums $Y_i = \sum_{j \le i} \tilde{r}_j$. Trend is Pearson correlation with position $i$. CV, Fano factor, and normalized successive difference use denominator $\bar{r} + 10^{-10}$ and are 0 when $\bar{r} \le 0$. Lag-1 autocorrelation and trend are 0 for denominators $\le 10^{-10}$. Skewness is 0 when $L < 3$ or $s \le 10^{-10}$; Hurst is 0.5 when $L < 8$ or $R$ or $s$ is $\le 10^{-10}$. References: \citep{ref53,ref54,ref55,ref65,ref66,ref108,ref109,ref110,ref113}.
\end{minipage}
\end{table}

\clearpage
\begin{table}[H]
\centering
\footnotesize
\captionsetup{position=top,skip=12pt}
\caption{Stress-test protocols and inference.}
\label{tab:si-protocol}
\begin{tabularx}{\linewidth}{@{}llXl@{}}
\toprule
\textbf{Category} & \textbf{Method} & \textbf{Description} & \textbf{Frontier} \\
\midrule
\multicolumn{4}{@{}l}{\textbf{Perturbation} (TOST equivalence at maximum intensity, $\delta_{\mathrm{EI}} = 0.08$, $\delta_{\mathrm{CI}} = 0.14$)} \\
\midrule
Network & Rewire & Degree-preserving edge swaps, fraction 0.05--0.35 (max 0.35) & Frozen \\
 & Densify & Uniform edge addition, fraction 0.05--0.35 (max 0.35) & Frozen \\
 & Sample & Uniform node removal, fraction 0.05--0.35 (max 0.35) & Frozen \\
\addlinespace
Invariants & Noise & Additive Gaussian noise, corpus-wide, 0.05--0.35 (max 0.35) & Frozen \\
 & Jitter & Multiplicative jitter, 0.05--0.35 (max 0.35) & Frozen \\
 & Subset & Random coordinate subset, retention 0.95--0.65 (max 0.65) & Frozen \\
\addlinespace
Process & Bootstrapping & Bootstrap to destroy temporal order 0.05--0.35 (max 0.05) & Frozen \\
 & Scaling & Power-law scaling, exponent 0.25--1.75 (max 0.25) & Frozen \\
 & Smoothing & Rolling-mean smoothing, window 3--15 intervals (max 15) & Frozen \\
\addlinespace
Signatures & Noise & Additive Gaussian noise, corpus-wide, 0.05--0.35 (max 0.35) & Frozen \\
 & Jitter & Multiplicative jitter, 0.05--0.35 (max 0.35) & Frozen \\
 & Subset & Random coordinate subset, retention 0.95--0.65 (max 0.65) & Frozen \\
\midrule
\multicolumn{4}{@{}l}{\textbf{Falsification} (one-sided inferiority to the original condition, Holm-adjusted)} \\
\midrule
Negative & Target Permutation & Cross-domain derangement of $y^{*}$ with $x$ and $z$ unchanged & Frozen + Retrain \\
 & Random Generation & Each feature redrawn uniformly within its observed corpus range & Frozen + Retrain \\
 & Vector Generation & Feature bootstrap plus Gaussian noise at pooled dispersion & Frozen + Retrain \\
\midrule
\multicolumn{4}{@{}l}{\textbf{Ablation} (one-sided non-inferiority of the log-additive decomposition, $\delta_{\mathrm{EI}} = 0.08$, $\delta_{\mathrm{CI}} = 0.14$)} \\
\midrule
Simple & Invariants Only & Frontier fitted on $x'$ alone & Retrain \\
 & Signatures Only & Frontier fitted on $z'$ alone & Retrain \\
\addlinespace
Complex & Joint Features & $x'$ and $z'$ concatenated into a single-stage fit & Retrain \\
 & Interaction Terms & Augmented with all 147 pairwise products & Retrain \\
\bottomrule
\end{tabularx}
\par\vspace{12pt}
\begin{minipage}{\linewidth}
\footnotesize
Stress tests use 30 model-fitting seeds (42--71), paired one-to-one with 30 independent stochastic perturbation realizations or applied to a single deterministic realization. Grids contain 7 equally spaced values between the listed endpoints; scaling includes the identity $\gamma = 1$. Feature interventions act on raw coordinates before fold-wise standardization. Additive-noise standard deviation is $s\sigma_j$, where $\sigma_j$ is the sample standard deviation of coordinate $j$ across all 25 original systems. Masking retains $\lfloor\rho p\rfloor$ coordinates, with $p = 21$ for invariants and $p = 7$ for signatures; noise, jitter draws, and masks are shared across systems within a realization. Network and process interventions recompute the affected feature vector while retaining original maximum targets. Complete bipartite node removal uses exact closed-form invariants; rewiring and densification use deterministic approximations (\hyperref[sec:si-note-bipartite]{Note~S1}), with rewiring interpolated toward configuration-model estimates \citep{ref57}. Rewiring and densification leave the complete graphs gwosc, seismic, and windmill unchanged. Falsified corpora are generated once with seed 42 as corpus-level negative controls. ``Max'' denotes the most severe setting, not necessarily the largest parameter: severity is intensity $\epsilon$ for network methods, intensity $s$ for noise and jitter, window $w$ for smoothing, $1 - \rho$ for retention $\rho$, $1 - s$ for bootstrapping, and $\lvert \ln \gamma \rvert$ for scaling exponent $\gamma$. Inference uses the 45 paired estimator-by-domain cells (9 estimators $\times$ 5 domains); metric aggregation, margins, and Holm-adjustment families are specified across test summaries (Tables~\ref{tab:si-perturb}--\ref{tab:si-ablation}).
\end{minipage}
\end{table}

\clearpage
\subsection*{Note S1: Complete Bipartite Perturbations}
\phantomsection
\label{sec:si-note-bipartite}
The 4 complete bipartite substrates (college, email, mooc, overflow) were represented by their fold sizes $m$ and $n$, with $N = m + n$ nodes and $M = mn$ edges, because their explicit graphs reach up to 7,200,000,000 edges. Node removal at intensity $\epsilon$ deleted $\lfloor \epsilon N \rfloor$ nodes drawn uniformly without replacement, counted the removed nodes in each fold, and evaluated the closed forms (Table~\ref{tab:si-features}) for the remaining complete bipartite graph. This computation is exact and used 30 realizations. Rewiring and densification used a deterministic degree-class approximation with 1 realization.

For degree classes $c$ with degree $d_c$ and $N_c$ nodes, and an edge count $M'$, let
\begin{equation*}
N = \sum_c N_c, \quad S_1 = \sum_c N_c d_c, \quad S_2 = \sum_c N_c d_c^{2}, \quad \bar{d} = \frac{S_1}{N}, \quad \kappa = \frac{S_2}{\max(S_1, 1)} - 1, \quad \overline{d^{-1}} = \frac{1}{N} \sum_c \frac{N_c}{\max(d_c, 1)}.
\end{equation*}
The approximation $\Phi$ assigned the following values and set articulation points, bridges, clustering, and assortativity to 0.
\begin{itemize}[leftmargin=1.5em,itemsep=0.2em]
\item Size: $|V| = N$, $|E| = M'$, and $d_{\max} = \max_c d_c$.
\item Cohesion: diameter $D = \max\{2, \ln \max(N, 2) / \ln \kappa\}$ when $\kappa > 1$ and $D = N$ otherwise, radius $D/2$, degeneracy $k^{*} = \min(d_{\max}, \sqrt{2M'})$, and $k$-core size $\sum_c N_c \mathbf{1}[d_c \ge k^{*}]$.
\item Degree statistics: variance $N^{-1} \sum_c N_c (d_c - \bar{d})^2$, bias-corrected skewness and excess kurtosis of the class-weighted degrees, degree entropy $-\sum_k p_k \ln p_k$ over degrees rounded to integers $k$ with vertex fractions $p_k$, and joint degree entropy $-2 \sum_k q_k \ln q_k$ with stub fractions $q_k \propto k p_k$.
\item Spectral moments: $\mathrm{tr}(L^2)/N = 1 + 2M'N / \max(S_1, 1)^2$, $\mathrm{tr}(L^3)/N = 1 + (S_2 - S_1)^3 / \max(N S_1^3, 1)$, $\mathrm{tr}(P^3)/N = 6 (S_2 - S_1)^2\, \overline{d^{-1}} / [\max(6 S_1, 1) \max(S_1, 1)]$, $\mathrm{tr}(P^4)/N = N^{-1} \sum_c N_c [d_c S_2 / \max(S_1^2, 1) + \max(\kappa, 0) / \max(d_c, 1)]$, and $\mathrm{tr}(A^4)/N = [S_2 + S_2^2 / \max(2 S_1, 1)] / N$.
\end{itemize}
Rewiring at intensity $\epsilon$ evaluated $\Phi$ on the classes of $K_{m,n}$, with degrees $(n, m)$, class sizes $(m, n)$, and $M' = M$, and interpolated each invariant as $(1 - \alpha) x_j + \alpha \Phi_j$ with $\alpha = 1 - (1 - \epsilon)^2$, where $x_j$ is the exact $K_{m,n}$ value. Densification added $a = \min(\lfloor \epsilon M \rfloor, C)$ edges among the $C = \binom{m}{2} + \binom{n}{2}$ within-fold vertex pairs, so that each pair was added with probability $q = a / C$. It evaluated $\Phi$ with degrees $(n + q(m - 1), m + q(n - 1))$, class sizes $(m, n)$, and $M' = M + a$. It then multiplied the $K_{m,n}$ articulation points, bridges, and assortativity by $1 - q$ and set clustering to $\min\{c_0 + (1 - c_0) q, 1\}$, where $c_0$ is the $K_{m,n}$ clustering. All other substrates were perturbed as explicit graphs.

\clearpage
\subsection*{Note S2: Corpus Processing Conventions}
\phantomsection
\label{sec:si-note-processing}
The \textbf{federal} series excludes awards with non-positive amounts or start dates outside 2011-01-01 to 2024-12-31. Zero speeds are treated as missing in \textbf{metrla} (8.1\% of readings) and as readings in \textbf{pemsbay} (0.003\%). \textbf{Montevideo} drops the first 4 hours removed by the loader lag and the final 20 hours of the 744-hour source that do not complete a day.

GPS timestamps in \textbf{gwosc} and \textbf{auger} were converted to UTC without leap-second correction; correction would move 19 \textbf{auger} events and no \textbf{gwosc} events to another day. The \textbf{crop} network retains 12 nodes and 30 edges from studies whose visits lack dates. Exact-name matching of \textbf{world} countries to lending metadata leaves 23 regional aggregates and spelling variants isolated.

\clearpage
\subsection*{SI References}
\phantomsection
\label{sec:si-sources}
\footnotesize
\begin{enumerate}[label=S\arabic*.,ref=S\arabic*,leftmargin=2.8em,itemsep=0.6em]
\item \label{src:S1} Vallisneri, M., Kanner, J., Williams, R., Weinstein, A. \& Stephens, B. The LIGO Open Science Center. \textit{J. Phys. Conf. Ser.} \textbf{610}, 012021 (2015). \url{https://doi.org/10.1088/1742-6596/610/1/012021}
\item \label{src:S2} U.S. Geological Survey. USGS Water Data for the Nation (1994). \url{https://doi.org/10.5066/F7P55KJN}
\item \label{src:S3} Pierre Auger Collaboration. Pierre Auger Observatory Open Data, version 3 [Dataset]. Zenodo (2024). \url{https://doi.org/10.5281/zenodo.10488964}
\item \label{src:S4} Albuquerque Seismological Laboratory/USGS. Global Seismograph Network (GSN - IRIS/USGS) [Dataset]. International Federation of Digital Seismograph Networks (1988). \url{https://doi.org/10.7914/SN/IU}
\item \label{src:S5} U.S. Geological Survey. Advanced National Seismic System (ANSS) Comprehensive Catalog (2017). \url{https://doi.org/10.5066/F7MS3QZH}
\item \label{src:S6} Meteostat. Meteostat weather and climate database. \url{https://meteostat.net} (accessed 1 September 2026).
\item \label{src:S7} Cassada, R. C. \& Russell, R. L. The dauerlarva, a post-embryonic developmental variant of the nematode \textit{Caenorhabditis elegans}. \textit{Dev. Biol.} \textbf{46}, 326--342 (1975). \url{https://doi.org/10.1016/0012-1606(75)90109-8}
\item \label{src:S8} Golden, J. W. \& Riddle, D. L. The \textit{Caenorhabditis elegans} dauer larva: developmental effects of pheromone, food, and temperature. \textit{Dev. Biol.} \textbf{102}, 368--378 (1984). \url{https://doi.org/10.1016/0012-1606(84)90201-X}
\item \label{src:S9} Allen-Perkins, A. et al. CropPol: A dynamic, open and global database on crop pollination. \textit{Ecology} \textbf{103}, e3614 (2022). \url{https://doi.org/10.1002/ecy.3614}
\item \label{src:S10} Hungarian Chickenpox Cases [Dataset]. UCI Machine Learning Repository (2021). \url{https://doi.org/10.24432/C5103B}
\item \label{src:S11} U.S. Food and Drug Administration. openFDA Drug Event API. \url{https://open.fda.gov/apis/drug/event/} (accessed 13 March 2026).
\item \label{src:S12} Guttag, J. CHB-MIT Scalp EEG Database, version 1.0.0 [Dataset]. PhysioNet (2010). \url{https://doi.org/10.13026/C2K01R}
\item \label{src:S13} U.S. Department of the Treasury. USAspending API. \url{https://api.usaspending.gov/} (accessed 25 August 2026).
\item \label{src:S14} Kumar, S., Zhang, X. \& Leskovec, J. Predicting dynamic embedding trajectory in temporal interaction networks. In \textit{Proc. 25th ACM SIGKDD Int. Conf. Knowledge Discovery and Data Mining} 1269--1278 (ACM, 2019). \url{https://doi.org/10.1145/3292500.3330895}
\item \label{src:S15} McAuley, J., Targett, C., Shi, Q. \& van den Hengel, A. Image-based recommendations on styles and substitutes. In \textit{Proc. 38th Int. ACM SIGIR Conf. Research and Development in Information Retrieval} 43--52 (ACM, 2015). \url{https://doi.org/10.1145/2766462.2767755}
\item \label{src:S16} Kumar, S., Spezzano, F., Subrahmanian, V. S. \& Faloutsos, C. Edge weight prediction in weighted signed networks. In \textit{Proc. IEEE 16th Int. Conf. Data Mining} 221--230 (IEEE, 2016). \url{https://doi.org/10.1109/ICDM.2016.0033}
\item \label{src:S17} World Bank. Projects \& Operations API and Country API. \url{https://search.worldbank.org/api/v2/projects} and \url{https://api.worldbank.org/v2/country} (accessed 29 September 2026).
\item \label{src:S18} Wikipedia Math Essentials [Dataset]. UCI Machine Learning Repository (2021). \url{https://doi.org/10.24432/C5V32F}
\item \label{src:S19} Paranjape, A., Benson, A. R. \& Leskovec, J. Motifs in temporal networks. In \textit{Proc. 10th ACM Int. Conf. Web Search and Data Mining} 601--610 (ACM, 2017). \url{https://doi.org/10.1145/3018661.3018731}
\item \label{src:S20} Panzarasa, P., Opsahl, T. \& Carley, K. M. Patterns and dynamics of users' behavior and interaction: network analysis of an online community. \textit{J. Am. Soc. Inf. Sci. Technol.} \textbf{60}, 911--932 (2009). \url{https://doi.org/10.1002/asi.21015}
\item \label{src:S21} OpenStreetMap contributors. OpenStreetMap [Dataset]. \url{https://www.openstreetmap.org} (accessed 10 November 2025).
\item \label{src:S22} Boeing, G. Modeling and analyzing urban networks and amenities with OSMnx. \textit{Geogr. Anal.} \textbf{57}, 567--577 (2025). \url{https://doi.org/10.1111/gean.70009}
\item \label{src:S23} Li, Y., Yu, R., Shahabi, C. \& Liu, Y. Diffusion convolutional recurrent neural network: data-driven traffic forecasting. In \textit{Int. Conf. Learning Representations} (2018).
\item \label{src:S24} Abbott, B. P. et al. GWTC-1: A Gravitational-Wave Transient Catalog of Compact Binary Mergers Observed by LIGO and Virgo during the First and Second Observing Runs. \textit{Phys. Rev. X} \textbf{9}, 031040 (2019). \url{https://doi.org/10.1103/PhysRevX.9.031040}
\item \label{src:S25} Abbott, R. et al. GWTC-2: Compact Binary Coalescences Observed by LIGO and Virgo during the First Half of the Third Observing Run. \textit{Phys. Rev. X} \textbf{11}, 021053 (2021). \url{https://doi.org/10.1103/PhysRevX.11.021053}
\item \label{src:S26} Abbott, R. et al. GWTC-2.1: Deep extended catalog of compact binary coalescences observed by LIGO and Virgo during the first half of the third observing run. \textit{Phys. Rev. D} \textbf{109}, 022001 (2024). \url{https://doi.org/10.1103/PhysRevD.109.022001}
\item \label{src:S27} Abbott, R. et al. GWTC-3: Compact Binary Coalescences Observed by LIGO and Virgo during the Second Part of the Third Observing Run. \textit{Phys. Rev. X} \textbf{13}, 041039 (2023). \url{https://doi.org/10.1103/PhysRevX.13.041039}
\item \label{src:S28} Government of Uruguay. National Open Data Catalog. \url{https://catalogodatos.gub.uy/} (accessed 1 September 2026).
\item \label{src:S29} Kunz, N. \& Gao, H. O. Global geolocated realtime data of interfleet urban transit bus idling. Preprint at \url{https://doi.org/10.48550/arXiv.2403.03489} (2024).
\item \label{src:S30} LIGO Scientific Collaboration, Virgo Collaboration \& KAGRA Collaboration. GWTC-4.0 Data Release Documentation [Dataset]. Gravitational Wave Open Science Center (2025). \url{https://doi.org/10.7935/aes8-px89}
\item \label{src:S31} Rader, R. et al. Non-bee insects are important contributors to global crop pollination. \textit{Proc. Natl Acad. Sci. USA} \textbf{113}, 146--151 (2016). \url{https://doi.org/10.1073/pnas.1517092112}
\end{enumerate}

\endgroup
\end{document}
